% GENERATED FILE. Edit canonical lessons, then run npm run generate:books.
% canonical-source-hash: fnv1a64:0dbc0c6821a58caa
% canonical-lessons: FR-C492-recompose, FR-C492-adopter, FR-C492-adoption, FR-C492-fratrie, FR-C492-famille-nombreuse

\chapter{Familles d'aujourd'hui}
\label{ch:fr-familles-d-aujourd-hui}
\hlchaptermodality{492}

\begin{chapteropening}
\textbf{By the end of this chapter:} \emph{I can describe different kinds of families.}

Complete the last lesson of chapter 492: I can describe different kinds of families.
\end{chapteropening}

\section[recomposé]{recomposé — blended (of a family)}

\label{lesson:FR-C492-recompose}



\begin{quote}
Before the new one: say the French for to break up, then the French for break-up.
\end{quote}

\subsection*{You'll want to know: recomposé}
\textbf{recomposé} — \textquotedblleft{}blended (of a family)\textquotedblright{}.

\begin{quote}
Nous sommes une famille recomposée avec cinq enfants.
\end{quote}

\begin{quote}
\emph{We are a blended family with five children.}
\end{quote}

\emph{Une famille recomposée} — parents and children from earlier unions living together.

\begin{cousinweb}
\emph{recomposé} is built from \emph{re-} and \emph{composé}, from \emph{composer} \textquotedblleft{}to put together\textquotedblright{}, which goes back to Latin \emph{compōnere}, reshaped after \emph{poser}; both rest on \emph{pōnere} \textquotedblleft{}to put\textquotedblright{} (Lewis \& Short).
\end{cousinweb}

\begin{grammarlens}[title={What you've built}]
One more word for telling, explaining and describing at B1.
\end{grammarlens}

\subsection*{Your turn}
\begin{itemize}
\raggedright
  \item \emph{Say it:} \emph{recomposé}
  \item \emph{Say it:} \emph{recomposé}, once more
  \item \emph{Say it:} a sentence of your own with \emph{recomposé}
  \item \emph{Recall:} say the French for to break up, then the French for break-up, then say \emph{recomposé} again
\end{itemize}

\begin{tcolorbox}[breakable,colback=teal!4,colframe=teal!35!black,title={Before you move on}]
What does recomposé mean? (\textquotedblleft{}Blended (of a family)\textquotedblright{}.) Say it once more.
\end{tcolorbox}
\section[adopter]{adopter — to adopt}

\label{lesson:FR-C492-adopter}



\begin{quote}
Before the new one: say the French for break-up, then the French for blended.
\end{quote}

\subsection*{You'll want to know: adopter}
\textbf{adopter} — \textquotedblleft{}to adopt\textquotedblright{}.

\begin{quote}
Mes voisins ont adopté une petite fille.
\end{quote}

\begin{quote}
\emph{My neighbours have adopted a little girl.}
\end{quote}

Also of an idea or a pet: \emph{adopter un chien}.

\begin{cousinweb}
\emph{adopter} is a learned borrowing from Latin \emph{adoptāre} \textquotedblleft{}to choose, to adopt\textquotedblright{} (EtymDB), built on \emph{optāre} \textquotedblleft{}to choose, to wish\textquotedblright{} (Lewis \& Short).
\end{cousinweb}

\begin{grammarlens}[title={What you've built}]
One more word for telling, explaining and describing at B1.
\end{grammarlens}

\subsection*{Your turn}
\begin{itemize}
\raggedright
  \item \emph{Say it:} \emph{adopter}
  \item \emph{Say it:} \emph{adopter}, once more
  \item \emph{Say it:} a sentence of your own with \emph{adopter}
  \item \emph{Recall:} say the French for break-up, then the French for blended, then say \emph{adopter} again
\end{itemize}

\begin{tcolorbox}[breakable,colback=teal!4,colframe=teal!35!black,title={Before you move on}]
What does adopter mean? (\textquotedblleft{}To adopt\textquotedblright{}.) Say it once more.
\end{tcolorbox}
\section[l'adoption]{l'adoption — adoption}

\label{lesson:FR-C492-adoption}



\begin{quote}
Before the new one: say the French for blended, then the French for to adopt.
\end{quote}

\subsection*{You'll want to know: l'adoption}
\textbf{l'adoption} — \textquotedblleft{}adoption\textquotedblright{}.

\begin{quote}
L'adoption a pris trois ans.
\end{quote}

\begin{quote}
\emph{The adoption took three years.}
\end{quote}

Feminine. \emph{Un enfant adopté} — an adopted child.

\begin{cousinweb}
\emph{l'adoption} is a learned borrowing from Latin \emph{adoptiō} \textquotedblleft{}adoption\textquotedblright{}, from \emph{adoptāre}, built on \emph{optāre} \textquotedblleft{}to choose\textquotedblright{} (Lewis \& Short).
\end{cousinweb}

\begin{grammarlens}[title={What you've built}]
One more word for telling, explaining and describing at B1.
\end{grammarlens}

\subsection*{Your turn}
\begin{itemize}
\raggedright
  \item \emph{Say it:} \emph{l'adoption}
  \item \emph{Say it:} \emph{l'adoption}, once more
  \item \emph{Say it:} a sentence of your own with \emph{l'adoption}
  \item \emph{Recall:} say the French for blended, then the French for to adopt, then say \emph{l'adoption} again
\end{itemize}

\begin{tcolorbox}[breakable,colback=teal!4,colframe=teal!35!black,title={Before you move on}]
What does l'adoption mean? (\textquotedblleft{}Adoption\textquotedblright{}.) Say it once more.
\end{tcolorbox}
\section[la fratrie]{la fratrie — siblings (as a group)}

\label{lesson:FR-C492-fratrie}



\begin{quote}
Before the new one: say the French for to adopt, then the French for adoption.
\end{quote}

\subsection*{You'll want to know: la fratrie}
\textbf{la fratrie} — \textquotedblleft{}siblings (as a group)\textquotedblright{}.

\begin{quote}
Dans ma fratrie, nous sommes quatre.
\end{quote}

\begin{quote}
\emph{There are four of us siblings.}
\end{quote}

A useful word with no plain everyday English match: all the brothers and sisters of one family.

\begin{cousinweb}
\emph{la fratrie} is a modern learned word built on Latin \emph{frāter} \textquotedblleft{}brother\textquotedblright{} (Lewis \& Short).
\end{cousinweb}

\begin{grammarlens}[title={What you've built}]
One more word for telling, explaining and describing at B1.
\end{grammarlens}

\subsection*{Your turn}
\begin{itemize}
\raggedright
  \item \emph{Say it:} \emph{la fratrie}
  \item \emph{Say it:} \emph{la fratrie}, once more
  \item \emph{Say it:} a sentence of your own with \emph{la fratrie}
  \item \emph{Recall:} say the French for to adopt, then the French for adoption, then say \emph{la fratrie} again
\end{itemize}

\begin{tcolorbox}[breakable,colback=teal!4,colframe=teal!35!black,title={Before you move on}]
What does la fratrie mean? (\textquotedblleft{}Siblings (as a group)\textquotedblright{}.) Say it once more.
\end{tcolorbox}
\section[la famille nombreuse]{la famille nombreuse — large family}

\label{lesson:FR-C492-famille-nombreuse}



\begin{quote}
Before the new one: say the French for adoption, then the French for siblings.
\end{quote}

\subsection*{You'll want to know: la famille nombreuse}
\textbf{la famille nombreuse} — \textquotedblleft{}large family\textquotedblright{}.

\begin{quote}
Elle vient d'une famille nombreuse de sept enfants.
\end{quote}

\begin{quote}
\emph{She comes from a large family of seven children.}
\end{quote}

In France, from three children on; such families get some discounts.

\begin{cousinweb}
\emph{la famille nombreuse} joins \emph{famille}, a learned borrowing from Latin \emph{familia} \textquotedblleft{}household\textquotedblright{}, and \emph{nombreux} \textquotedblleft{}numerous\textquotedblright{}, from Latin \emph{numerōsus}, built on \emph{numerus} \textquotedblleft{}number\textquotedblright{} (Lewis \& Short).
\end{cousinweb}

\begin{grammarlens}[title={What you've built}]
One more word for telling, explaining and describing at B1.
\end{grammarlens}

\subsection*{Your turn}
\begin{itemize}
\raggedright
  \item \emph{Say it:} \emph{la famille nombreuse}
  \item \emph{Say it:} \emph{la famille nombreuse}, once more
  \item \emph{Say it:} a sentence of your own with \emph{la famille nombreuse}
  \item \emph{Recall:} say the French for adoption, then the French for siblings, then say \emph{la famille nombreuse} again
\end{itemize}

\begin{tcolorbox}[breakable,colback=teal!4,colframe=teal!35!black,title={Before you move on}]
What does la famille nombreuse mean? (\textquotedblleft{}Large family\textquotedblright{}.) Say it once more.
\end{tcolorbox}
